\documentclass[12pt]{article}
\usepackage[T1]{fontenc}
\usepackage[utf8]{inputenc}
\usepackage{lmodern}
\usepackage{textcomp}
\usepackage{enumitem}
\usepackage{geometry}
\usepackage{graphicx}
\usepackage{amsmath, amssymb}
\geometry{margin=1in}
\begin{document}
\begin{center}
Sociology - Graduate Level Assessment 15 \
Total Marks: 40 \
Answer all questions.
\end{center}

\section*{Multiple Choice (10 marks)}
\begin{enumerate}[label=\arabic*.]
\item Bernstein thought that using 'restricted codes' of language disadvantaged pupils at school because:
\begin{enumerate}[label=(\alph*)]
    \item this pattern of speech made them the target of bullying
    \item they referred to explicit, context independent meanings
    \item they prevented children from communicating outside of their peer groups
    \item they involved short, simple sentences with a small vocabulary
\end{enumerate}
\item World-affirming religions:
\begin{enumerate}[label=(\alph*)]
    \item embrace conventional cultural values but offer new means of achieving them
    \item react against the loss of any meaningful religious content in the teachings of churches
    \item adopt an attitude of mild disapproval towards mainstream social values
    \item reject both the goals and means of conventional society and provide utopian alternatives
\end{enumerate}
\item A survey should avoid asking:
\begin{enumerate}[label=(\alph*)]
    \item fixed-choice questions
    \item short questions
    \item leading questions
    \item funnelled questions
\end{enumerate}
\item The 'correspondence principle' (Bowles \& Gintis) suggests that:
\begin{enumerate}[label=(\alph*)]
    \item schools prepare children for work by teaching them to be obedient
    \item teachers and parents tend to have similar attitudes to learning
    \item children who write lots of letters develop a better grasp of language
    \item boys' and girls' educational achievements have recently become similar
\end{enumerate}
\item Herberg's (1955) study of religion in America suggested that:
\begin{enumerate}[label=(\alph*)]
    \item ethnic minorities practised religion to achieve social acceptance in the culture
    \item mainstream faiths were becoming increasingly identified with national identity
    \item the moral teachings of the main religions were becoming relatively similar
    \item all of the above
\end{enumerate}
\end{enumerate}

\section*{True / False (10 marks)}
\textit{Answer True or False}
\begin{enumerate}[label=\arabic*.]
\setcounter{enumi}{5}
\item True or False: The Mafia is primarily considered a type of white collar crime, focusing on non-violent financial fraud.
\item True or False: The nineteenth century saw a significant rise in women occupying high-status positions in industrial production settings.
\item True or False: Human behavior is meaningful and varies significantly between individuals and cultures, making societal study distinct from natural sciences.
\item True or False: A moral panic occurs when the media exaggerate reports of deviant groups, leading to heightened public anxiety and hostility.
\item True or False: Murray argued that the 'underclass' were primarily individuals who chose not to participate in the labor market due to personal preference.
\end{enumerate}

\section*{Long Form Response (20 marks)}
\textit{Answer each question with a clear and complete explanation.}
\begin{enumerate}[label=\arabic*.]
\setcounter{enumi}{10}
\item Discuss the concept of scapegoating in sociology, analyzing its implications for social cohesion and conflict within a society.
\item Discuss the functionalist perspective in sociology, analyzing how it explains the mechanisms through which society maintains order and stability in various social systems.
\end{enumerate}

\vfill
\noindent\textit{End of Paper}
\end{document}